% !PRINTED 177
\weberpage{177}
Die Funktion $\psi(\omega)$ kann möglicherweise bei gewissen Transformationen
 $S$ ungeändert bleiben (z.\,B. immer bei der identischen
 Transformation). Alle Transformationen, die eine Funktion
 $\psi(\omega)$ ungeändert lassen, bilden eine \gesperrt{Gruppe} $\mathfrak{G}$. Denn ist
\[
\psi\,|\,S = \psi, \qquad \psi\,|\,S_1 = \psi,
\]
so ist auch
\[
\psi\,|\,SS_1 = \psi\,|\,S_1 = \psi.
\]
Die Gruppe $\mathfrak{G}$ ist ein Teiler der Gruppe $\mathfrak{L}$ aller linearen
Transformationen.

Ist $\mathfrak{G}$ ein Teiler von $\mathfrak{L}$ von \gesperrt{endlichem Index $(\mathfrak{L}, \mathfrak{G})$, und
$\psi(\omega)$ eine Funktion von $\omega$ von der Eigenschaft, daß auch
umgekehrt}
\[
\psi(\omega') = \psi(\omega)
\]
\gesperrt{einen linearen Zusammenhang}
\[
\omega' = \frac{\gamma+\delta\omega}{\alpha+\beta\omega}
\]
\gesperrt{zur Folge hat, in dem $\begin{pmatrix} \alpha, & \beta \\ \gamma, & \delta \end{pmatrix}$ eine zu $\mathfrak{G}$ gehörige Transformation
 ist, so heißt $\psi(\omega)$ eine zur Gruppe $\mathfrak{G}$ gehörige}
Modulfunktion.

Wir beschränken uns auf die Betrachtung solcher Modulfunktionen,
die mit dem Modul $\varkappa^2$, also auch mit $j(\omega)$, in einem
\gesperrt{algebraischen Zusammenhange} stehen.

Unter dieser Voraussetzung läßt sich der folgende allgemeine
Satz aussprechen:

\gesperrt{Wenn $\psi(\omega)$ zur Gruppe $\mathfrak{G}$ gehört und $\chi(\omega)$ durch die
Transformationen von $\mathfrak{G}$ ungeändert bleibt, so ist $\chi(\omega)$}
\textbf{rational} \gesperrt{durch $\psi(\omega)$ ausdrückbar.}

Denn zunächst ist $\chi(\omega)$ eine algebraische Funktion von $\psi(\omega)$,
und $\psi(\omega)$ kann als algebraische Funktion von $j(\omega)$ jeden Wert
annehmen. Zu jedem Wert von $\psi(\omega)$ gehört aber nach \gesperrt{Voraussetzung
 nur ein Wert von $\chi(\omega)$, und daher ist $\chi(\omega)$ als einwertige
 algebraische Funktion von $\psi(\omega)$ rational.}

Das Studium der in $\mathfrak{L}$ enthaltenen Gruppen und der zu ihnen
gehörigen Modulfunktionen bildet den Hauptgegenstand des großen
Werkes von \gesperrt{Klein und Fricke}: \glqq Vorlesungen über die Theorie
der elliptischen Modulfunktionen\grqq{} (2 Bde., Leipzig, Teubner,
1890, 1892). Wir betrachten hier nur einige ganz spezielle Fälle
dieser Gruppen und Funktionen, auf die wir im Verlauf unserer
Untersuchungen gestoßen sind.
